\section{课程导论：这门课在讲什么}

这节课是 Berkeley ``Large Language Model Agents'' 的开场导论，由 Dawn Song 介绍课程定位、教学团队和学习目标。她把 Agent 放在 LLM 快速进化的大背景下讨论：语言模型已经足够强，但若想解决真实任务，就必须让模型与环境、工具、记忆和反馈形成闭环。

\begin{figure}[H]
\centering
\includegraphics[page=2,width=0.9\textwidth]{slides.pdf}
\caption{课程教学团队与协作角色。导论首先明确师资、助教和课程支持资源。}
\end{figure}

\begin{knowledgebox}{这门课的出发点}
课程并不把 Agent 视为单次 prompt 的技巧集合，而是把它定义为一个完整系统：LLM 负责推理与规划，外部模块负责工具使用、记忆、检索和环境交互。只有把这些能力串起来，LLM 才能从``回答问题''转向``完成任务''。
\end{knowledgebox}

\subsection{大模型为什么自然走向 Agent}

导论指出，近两年 LLM 的能力提升极快，但纯文本输入到纯文本输出的模式很快触到了上限。真实世界任务通常不是一次性问答，而是需要：
\begin{itemize}
    \item 与环境交互，读取页面、文件、数据库或设备状态；
    \item 保留任务上下文和长期记忆；
    \item 在失败后重试、调整计划，并从反馈中修正行为；
    \item 调用外部工具弥补模型在计算、检索和执行上的短板。
\end{itemize}

\begin{importantbox}{从 LLM 到 Agent 的关键跳变}
一旦模型拥有了记忆、工具使用、检索和反馈回路，它的行为单位就不再是``回答一句话''，而是``在多步交互中推进任务''。这也是 Agent 研究和普通 chatbot 应用的根本分界线。
\end{importantbox}

\subsection{本章小结}

课程导论先把问题讲清楚了：Agent 不是给 LLM 再包一层 UI，而是把模型放进一个能持续观察、规划、行动和修正的执行框架中。

